%% ============================================================
%%  bookedge-demo.tex —— 「书边索引」模板示例书稿
%%  编译：xelatex bookedge-demo.tex（跑两遍以上，书边索引需要 aux 信息）
%% ============================================================
\documentclass{bookedge}

\bookedgetitle{书的模样}
\title{书的模样}
\epigraph{书刊版式设计的节奏与秩序}
\author{刘小洋\quad 编著}
\publishinfo{2026 年 11 月\quad 第 1 版\quad ·\quad 示例书稿，仅用于模板演示}

\begin{document}

% ------------------------------------------------------------
\maketitle                       % 封面（书口十色彩条自动绘制）

% 扉页
\thispagestyle{empty}
\begin{titlepage}
  \null\vspace*{92mm}
  \begin{center}
    {\sffamily\bfseries\fontsize{20}{26}\selectfont\color{black!85}书的模样\par}
    \vskip 5mm
    {\sffamily\fontsize{10}{14}\selectfont\color{black!50}书刊版式设计的节奏与秩序\par}
  \end{center}
\end{titlepage}

% 版权页
\thispagestyle{empty}
\begingroup
  \vspace*{158mm}
  \setlength{\baselineskip}{16pt}
  \noindent{\sffamily\footnotesize\color{black!55}
    书的模样——书刊版式设计的节奏与秩序\\
    刘小洋\quad 编著\\
    示例书稿 · 由 bookedge 模板排版生成\\
    开本\quad 170mm $\times$ 240mm\quad 1/16\\
    2026 年 11 月第 1 版\par}
\endgroup
\clearpage

% ------------------------------------------------------------
\chapter*{前言}
\addcontentsline{toc}{chapter}{前言}

拿起新书之前，很多读者的手指会先在书口上摩挲一遍——那些凸起的、按颜色排列的小色块，就是全书的地图。字典用铝箔或染色标签把数千个词条切成几十个便于翻检的区段；工具书用拇指索引（thumb index）把每一章钉在书口上。这本书的模板把同样的思路带进了 \LaTeX{} 排版：每一章拥有一种专属颜色，书口上对应位置的色块被点亮，翻动书页时，你随时知道自己身处全书哪一段旅程。

版式设计的目标从来不是炫技，而是降低读者的认知成本。书边索引是成本最低、收益最高的导航装置之一：它不占用版心一寸空间，却在每一次翻页时完成一次定位。本书即以自身为范例——你可以边读内容，边体会这套版式系统如何工作。

\bigskip
\noindent{\sffamily\color{black!60}刘小洋\quad 2026 年秋}

\chapter*{使用说明}
\addcontentsline{toc}{chapter}{使用说明}

\begin{entry}{拇指索引}{Thumb Index}
在书口（外缘）按章节等分排列的一列色块。当前章的色块以实色点亮并加宽，其余章节显示为浅色底纹，一眼即可定位当前进度与全书结构。
\end{entry}

\medskip
本书的导航体系由三层构成，自外向内分别是：

\begin{itemize}
  \item \textbf{书边索引}：位于书口，全书常驻。单数页在右缘，双数页在左缘（与装订后的纸边一致）；
  \item \textbf{章首导航尺}：位于每章首页顶部，以缩略色块预告全书结构，当前章高亮；
  \item \textbf{目录色签}：目录中每章条目前的小色签与书边索引同色，建立「目录—书口」的颜色映射。
\end{itemize}

\bigskip
\noindent 下图示意了书边索引在成书上的位置与构成：

\begin{center}
\begin{tikzpicture}[x=1mm, y=1mm]
  % 书页示意（单数页，书口在右）
  \draw[black!25, line width=0.5pt, fill=black!2] (0,0) rectangle (62,88);
  \draw[black!12, line width=0.4pt] (11,10) rectangle (44,78);
  \node[font=\sffamily\fontsize{6}{7}\selectfont\color{black!45}] at (27.5,74) {版心};
  \node[font=\sffamily\fontsize{6}{7}\selectfont\color{black!45}] at (7.2,6.5) {订口};
  % 书边色块（示意 6 章）
  \foreach \i/\c in {1/beC1,2/beC2,3/beC3,4/beC4,5/beC5,6/beC6}{%
    \pgfmathsetmacro\yy{78-(\i-1)*13}
    \ifnum\i=3
      \fill[rounded corners=1.2pt, \c] (62,\yy-12) rectangle (70,\yy-0.8);
      \node[rotate=90] at (66.2,{\yy-6.4}) {\textcolor{white}{\sffamily\bfseries\fontsize{6.5}{6.5}\selectfont 0\i}};
      \draw[->, \c!80!black, line width=0.5pt] (78,\yy-6) -- (70.8,\yy-6);
      \node[anchor=west, font=\sffamily\fontsize{6.5}{8}\selectfont\color{\c!70!black}] at (78.5,\yy-6) {当前章：实色加宽};
    \else
      \fill[rounded corners=1.2pt, \c!14] (62,\yy-12) rectangle (69,\yy-0.8);
      \node[rotate=90] at (65.6,{\yy-6.4}) {\textcolor{\c!55}{\sffamily\fontsize{5.5}{5.5}\selectfont 0\i}};
    \fi
  }
  \draw[->, black!50, line width=0.5pt] (52,86) -- (61,82);
  \node[anchor=east, font=\sffamily\fontsize{6.5}{8}\selectfont\color{black!55}] at (51,87) {书口（外缘）};
  \node[font=\sffamily\fontsize{6}{7}\selectfont\color{black!40}] at (31,-4) {单数页：索引在右缘\quad 双数页：索引在左缘};
\end{tikzpicture}
\end{center}

\bigskip
\noindent 三点使用约定：

\begin{enumerate}
  \item 全书章数 $\le 10$ 时效果最佳；超过 10 章时色板自动循环，索引条会变薄，建议改以「部」分色；
  \item 目录之前的页面（封面、扉页、版权页）不显示索引；从目录页起，索引以「全浅色」状态出现，相当于全书的图例；
  \item 每章的专属色同时用于章号、节号、要点框、目录色签与页码徽章，形成章内一致的视觉氛围。
\end{enumerate}

% 从目录页起启用书边索引
\bookedgetabson
\tableofcontents

% ------------------------------------------------------------
\mainmatter

\part{基础原理}

% ============================================================
\chapter{版式的骨骼：网格系统}

拿到一份书稿，设计师做的第一件事不是选字体，也不是挑颜色，而是画格子。网格（grid）是版面的隐形骨架：它规定了文字从哪里开始、在哪里停住、图片如何对齐、空白出现在何处。读者看不见网格，却能感觉到它——「这本书读起来很舒服」，多半是网格在起作用。

\section{为什么需要网格}

没有网格的版面，就像没有梁柱的房子：每一页都要临场决定边界，十页之后必然失控。网格的价值在于把决定权交给一个事先想清楚的系统：

\begin{itemize}
  \item \textbf{一致性}：所有页面的版心、栏宽、行位彼此对齐，翻页时视线不需要重新适应；
  \item \textbf{效率}：图文排入预先划定的格位，不必逐页斟酌；
  \item \textbf{灵活性}：网格不是牢笼，通栏、跨栏、出血图都是在网格内做出的例外，例外因系统的存在而显得从容。
\end{itemize}

\begin{entry}{网格系统}{Grid System}
把版面按一定模数划分为若干栏与行，并以此约束文字、图片与空白位置的版式方法。源于 20 世纪中期的瑞士平面设计，代表人物为约瑟夫·米勒–布罗克曼（Josef Müller-Brockmann），其《平面设计中的网格系统》一书对此有系统论述。
\end{entry}

\section{本书的网格}

本书版心置于 170mm $\times$ 240mm 的开本上，内白边 21mm、外白边 31mm，版心宽约 118mm，可容纳三十二个五号字。外缘特意留出的宽白边并非浪费：它是书边索引的容身之所，也是翻阅时拇指按压的地方——功能性的留白。

\begin{notebox}{要点}
网格的模数应由正文推出，而不是先画格子再塞文字。先确定字号与每行字数，再决定栏宽与版心，最后才轮到页边距的分配——顺序反了，字会迁就格子，读起来处处别扭。
\end{notebox}

行距是网格的时间维度。本书正文行距为字号的 1.42 倍，行间留有恰好能「呼吸」而不至于松散的空隙。行距一旦确定，标题、图注、框文的间距都应与之成倍数关系，版面才会产生节奏。

% ============================================================
\chapter{字体的声音：中西文搭配}

如果把版面比作房间，字体就是家具：同样大小的房间，摆什么家具、怎么摆，决定了它是书房还是客厅。中文书籍的字体搭配还多一层难度——要在汉字与拉丁字母两套风格迥异的系统之间求得统一。

\section{衬线与无衬线的分工}

正文字体要「隐形」：读者应当忘记字的存在，只记得内容。宋体（衬线体）横细竖粗、字形端庄，适合长文阅读；黑体（无衬线体）笔画均匀、视觉响亮，适合标题与导引。本书的分工正是如此：

\begin{itemize}
  \item 正文：宋体，五号，行距 1.42 倍——安静、连续、可长时间阅读；
  \item 标题与导引：思源黑体——从 50pt 的章号到 5.5pt 的导航尺，一族字重覆盖全部层级；
  \item 数字与西文：与正文宋体气质相近的衬线西文，避免西文「浮」在汉字之上。
\end{itemize}

\begin{entry}{字重}{Weight}
同一字族的不同笔画粗细等级，如 Light、Regular、Medium、Bold。设计系统的层级感主要来自字号与字重的配合，而非滥用加粗——全页加粗等于没有加粗。
\end{entry}

\section{层级的表情}

标题不必大，但要有表情。本书章节页以超大号数字配小号「CHAPTER」字样，大与小、实与虚的对比制造出翻页时的仪式感；正文内的节标题只用 13pt，靠颜色与字重区分层级。层级是靠关系表达的，不是靠尺寸堆叠。

\begin{notebox}{要点}
一页之内字族不宜超过三种：正文一族、标题一族、标注一族。每多引入一个字族，版面的「口音」就重一分，读者的注意力就被多分走一分。
\end{notebox}

% ============================================================
\chapter{色彩的秩序：从调色板到页边}

颜色是这本书的主角。十条低饱和色的色带排在书口上，像一排安静的哨兵：朱砂、赭橙、金棕、苔绿、松青、靛蓝、紫藤、玫紫、蓝灰、棕褐——每种颜色对应一章，从封面一直贯穿到最后一页。

\section{功能性用色}

本书的用色原则是：颜色只承担信息，不做装饰。每一种颜色都回答一个问题——

\begin{itemize}
  \item 我在哪一章？——书边索引与章首导航尺以实色回答；
  \item 这是哪个层级？——节号、列表符号以章色回答；
  \item 这是什么性质的框？——词条、要点、示例三种框各有固定色相。
\end{itemize}

选择色板时遵循三条约束。其一，全部取低饱和色，长时间阅读不刺眼，批量印刷不偏色；其二，饱和度够深的颜色才配白字，保证小色块上的章号可读；其三，相邻色相拉开距离，书口上一排色块依次可辨。

\begin{entry}{拇指索引}{Thumb Index}
在书口按内容分区排列的可辨识标记，词典式导航的传统形式。印刷上多以专色或染色标签实现，数字排版中则以版面外围的色块模拟。
\end{entry}

\section{颜色的克制}

一个常见的误区是把颜色当作装饰贴纸。当每一页都五颜六色时，颜色就失去了指路的能力。本书的设计是让颜色安静地工作：正文字永远保持墨色，色只出现在「路标」上——章号、节号、框线、页码徽章与书口。翻到任何一页，颜色的信息密度恰到好处：你看它一眼就知道身处何处，不看它也不受打扰。

\begin{notebox}{要点}
检验用色是否克制的方法很简单：把页面转成灰度。如果灰度版仍然层级分明、导航清晰，说明颜色用得有纪律；如果灰度版一片混沌，说明你是在用颜色遮掩层级设计的懒惰。
\end{notebox}

% ============================================================
\chapter{页边的风景：版心与留白}

页边距是版面中最容易被低估的元素。新手总想把字排满，仿佛白边是浪费；老手则知道，书页的呼吸感全部来自边缘。

\section{四条边的不同使命}

装订边的白边要承受翻开时的张力，不能太窄；外缘的白边要容纳翻阅的拇指与本书的书边索引，所以最宽；上边的白边给页眉与页码徽章留出位置；下边的白边托住版心，给脚注与归宿感。本书四边的比例约为 $21:31:24:28$（内：外：上：下，单位 mm），外缘明显宽于订口——这是书口索引方案的必然选择。

\begin{examplebox}{示例：双页跨页的节奏}
摊开一本书，视线先落在跨页中缝附近，再向两侧展开。因此订口可以略窄，让两个版心在视觉上彼此靠近；外口则宽出数毫米，成为整幅跨页的「画框」。索引色块正位于画框的边缘，既是装饰，也是坐标。
\end{examplebox}

\section{留白不是空}

留白是版式中最便宜的奢侈。段与段之间的空、图与文之间的空、章节页下部的大片空——它们不是「没排完」，而是刻意留下的停顿。章首页底部约三分之一版面留空，是全书节奏中最大的呼吸：翻过信息密度很高的一章正文，抵达新章的首页，视线在留白处稍作休息，再开始下一段旅程。

\begin{notebox}{要点}
页边的宽窄是一个系统决定，不是逐页的局部决定。一旦确定，所有页面共用同一组数字——读者对版面的信任感，正来自这种不被察觉的规律性。
\end{notebox}

\part{实践进阶}

% ============================================================
\chapter{导航的设计：目录、书眉与书边索引}

导航系统回答读者永恒的三个问题：我在哪里？我能去哪里？怎么回去？本书的答案是三层嵌套的色块系统——书边索引负责「我在哪里」，目录色签负责「我能去哪里」，章首导航尺负责两者之间的换算。

\section{三层导航的分工}

\begin{itemize}
  \item \textbf{书边索引（全书常驻）}：实色块的位置即当前章。不用翻到目录，拇指指腹在书口一摸便知进度；
  \item \textbf{章首导航尺（章首页）}：一排小色块按章序排列，当前章实色。它同时预告「这一章前后还有什么」，让章节不再是一串孤立的名字；
  \item \textbf{目录色签（目录页）}：每章条目前的彩色编号签与书口同色。从目录查到目标章，再到书口找到对应颜色，两次颜色匹配完成导航闭环。
\end{itemize}

\begin{entry}{书眉}{Running Head}
版心上方的常驻信息行，通常承印书名与章名。本书页眉取极简方案：内侧放章名或书名，外侧放一枚章色页码徽章，不加横线——颜色本身就是分隔。
\end{entry}

\section{导航即总结}

好的导航系统还有一个副产品：它是全书的抽象画。只看书口，你看到的是全书的等分切分与当前进度；只看导航尺，你看到的是各章的颜色序列。**导航设计的第一性原理，是让结构可见。** 读者未必意识到自己在读色块，但色块一直在替他们记路。

\begin{notebox}{要点}
导航元素应当「可预期」：每一页的书口、页眉、页码出现在完全相同的位置。变化等于噪音；常量才能被依赖。
\end{notebox}

% ============================================================
\chapter{图文的呼吸：插图与表格}

图与表是版面的变奏。正文是稳定的行进，图表是即兴的华彩——但华彩不能跑调，必须落回网格。

\section{图的宽度语言}

插图宽度应表达「信息量级」：半栏图是小注脚，通栏图是要点，跨页出血图是宣言。宽度即语气，全用通栏图等于全程大喊。本书示例的三种框——词条、要点、示例——宽度均与版心一致，仅以左侧色条与底色深浅区分语气，不打断行的节奏。

\begin{examplebox}{示例：图文距}
图与上文的关系必须近于图与下文的关系，读者才不会把图误挂到下一段。这一原则（亲和原则）同样适用于标题与节文、图注与图。
\end{examplebox}

\section{表格的安静}

表格的美德是低调。线越少越好：只保留表头下的横线与收尾横线，竖线全部省去；数字右对齐，文字左对齐；表头用无衬线体，表身用正文字体。一张安静表格的信息，比一张满线表格的信息更容易被读进去。

\begin{entry}{亲和原则}{Proximity Principle}
格式塔心理学的基本组织原则之一：距离近的元素被知觉为一个整体。版式上体现为「关系近的元素距离近，关系远的元素距离远」。
\end{entry}

% ============================================================
\chapter{节奏与重复：系列书的连续性}

一本书的美是运气，一套书的美是系统。当读者把第二册放到第一册旁边时，他们期待的不只是相似，而是「同一件事的延续」。

\section{可复用的骨架}

系列书的连续性来自三个可复用层：几何层（开本、版心、页边比例完全一致）、色彩层（同一色板，各册按序取色或按主题取色）、构件层（框式、页眉、页码徽章、章首页结构全部复用）。本书模板正是把这三层固化为一个类文件：换一部书稿，改标题与文字，系统不动。

\begin{examplebox}{示例：色板的接力}
设想本系列共三册。第一册用色板的前十个颜色；第二册沿用同一色板但整体加深一档；第三册回到原始色板但交换冷暖区段。三册放在一起，读者一眼认出「一家人」，又能分辨长幼。
\end{examplebox}

\section{重复中的变化}

纯粹的重复令人厌倦，纯粹的变化令人迷失。系列设计的分寸感在于：骨架不变，表情微调。章首页的大数字可以每册换一种排法，书口的色带可以每册换一个渐变方向——骨架承载「认得出」，表情承载「有新意」。

\begin{notebox}{要点}
给系列立规矩最好的时机是第一册开工之前，而不是第三册返工的时候。把规矩写成模板（代码或样式表），比写成文件更有约束力。
\end{notebox}

% ============================================================
\chapter{印刷的落地：从屏幕到纸张}

屏幕上的完美版式，到纸上一落印就可能走样。最后这一章谈三件落地的事：颜色、纸张与装订——它们决定书边索引这套方案最终能否兑现。

\section{颜色：从 RGB 到 CMYK}

屏幕色（RGB）与印刷色（CMYK）的色域不同，屏幕上鲜亮的颜色到了纸上往往发闷。本书色板特意选择低饱和色，就是为了缩小两端差距。批量印制前应打样核对，尤其是书口色带这类大面积色块——轻微的偏色在整条色带上会被放大成明显的色差。

\section{装订与索引的物理性}

本书的书边索引以印刷色块模拟拇指索引。若要做成真正的实体索引，可在印后加工中沿色块边缘模切（die-cut）出凸耳，或在书口刷色边。无论哪种工艺，都请与印厂确认：

\begin{itemize}
  \item 模切凸耳需要每章一个刀模，章数即成本，十章以内经济；
  \item 刷边（edge coloring）按书口整体计价，与章数无关，但只能呈现渐变或多段色，不能带编号；
  \item 锁线装订的书口最平整，无线胶订的书口在翻阅后会有轻微毛边，索引色块宜留出 2mm 容差。
\end{itemize}

\begin{entry}{出血}{Bleed}
图文延伸到裁切线之外的部分，裁切时被切掉，用以避免裁切误差造成的露白。书口满边色块在印刷时即按出血处理。
\end{entry}

\section{交付清单}

把模板交出去之前，最后核对一遍：两遍编译后书口索引的章数与实际章数一致；目录色签颜色与书口对应章颜色逐一对上；单双页索引分别位于右、左缘；灰度打印一版检查层级。这四项通过，这本书就可以下厂了。

\bigskip
\begin{notebox}{要点}
排版设计的终点不是屏幕上的 PDF，而是读者手中翻开的那一叠纸。所有的设计决定，最终都要在纸张、油墨与装订线上接受检验。
\end{notebox}

% ============================================================
\chapter*{附录：五分钟上手本模板}
\addcontentsline{toc}{chapter}{附录：五分钟上手本模板}

\begin{enumerate}
  \item 复制 \texttt{bookedge.cls} 与你的正文文件放在同一目录，正文文件以 \texttt{\string\documentclass\{bookedge\}} 开头；
  \item 用 \texttt{\string\title}、\texttt{\string\author}、\texttt{\string\epigraph}、\texttt{\string\publishinfo} 填写封面信息；
  \item 在 \texttt{\string\tableofcontents} 之前调用 \texttt{\string\bookedgetabson} 启用书边索引；
  \item 正文照常使用 \texttt{\string\chapter}、\texttt{\string\section}，词条用 \texttt{entry} 环境、要点用 \texttt{notebox}、示例用 \texttt{examplebox}；
  \item 用 \texttt{xelatex} 编译两遍以上（第二遍起总章数写入辅助文件，书口索引自动对位）。
\end{enumerate}

\bigskip
\noindent 想换色板？在 \texttt{bookedge.cls} 的「色板」一节修改十条 \texttt{\string\definecolor} 即可，全书与书口索引将同步更新。

\end{document}
